% !TeX root = ../informe.tex
%=======================================================================
%  config/plantilla.tex — mecanismos propios de la plantilla
%
%  Dos interruptores gobiernan qué se ve. ANTES DE ENTREGAR pon los dos en
%  «false»: desaparecen las cajas de instrucción y los marcadores quedan en
%  negro, sin que haya que borrar nada del texto.
%=======================================================================

%------------------------ Interruptores --------------------------------
\newif\ifguia          \guiatrue        % \guiafalse  oculta las instrucciones
\newif\ifmarcadores    \marcadorestrue  % \marcadoresfalse deja los [Insertar ...] en negro

%------------------------ Cajas de instrucción -------------------------
\definecolor{guiaback}{HTML}{F3F4F6}
\definecolor{guiaframe}{HTML}{9AA3AF}
\definecolor{marcador}{HTML}{B45309}

% \instruccion{...} — qué va en cada sitio. Se apagan con \guiafalse.
\newcommand{\instruccion}[1]{%
  \ifguia
    \begin{tcolorbox}[enhanced,breakable,
        colback=guiaback,colframe=guiaframe,boxrule=0.4pt,arc=1pt,
        left=6pt,right=6pt,top=4pt,bottom=4pt,
        fontupper=\small\itshape]
      \textbf{\textup{Instrucción.}}\enspace #1
    \end{tcolorbox}%
  \fi}

% \ph{...} — marcador de relleno, el equivalente de «[Insertar X]».
% En ámbar mientras se escribe; en negro al entregar, para que salte a la
% vista si alguno se quedó sin rellenar.
\newcommand{\ph}[1]{%
  \ifmarcadores{\color{marcador}\textsf{[#1]}}\else\textsf{[#1]}\fi}

%--------------------------- Notas de pendiente ------------------------
\setlength{\marginparwidth}{2.4cm}
% Quita «disable» para volver a ver las notas de pendiente. Con la opción
% puesta, no queda ninguna en el PDF aunque sigan escritas en el fuente.
\usepackage[spanish,colorinlistoftodos,prependcaption,
            textsize=footnotesize]{todonotes}

% \pend{...}   nota al margen
% \pendin{...} nota en línea, para textos largos o dentro de tablas
% \falta{...}  pendiente que bloquea la entrega
\newcommand{\pend}[1]{\todo[color=orange!35]{#1}}
\newcommand{\pendin}[1]{\todo[inline,color=orange!25]{#1}}
\newcommand{\falta}[1]{\todo[color=BrickRed!35]{#1}}

%----------------------- Identificadores trazables ---------------------
% Se usan en lugar de escribir el identificador a mano: crean hiperenlace
% y mantienen honestas las referencias cruzadas entre secciones. Si el
% destino no existe, LaTeX avisa y en el PDF sale «??» — que es justo lo
% que se quiere: un identificador inventado no pasa desapercibido.
%
%   \rn{02}      RN-02   regla de negocio        (sección 6)
%   \adrref{01}  ADR-01  decisión de arquitectura (sección 4)
%   \vista{02}   V-02    vista arquitectónica     (sección 4)
%   \crit{3}     CA-3    criterio de aceptación   (sección 7)
%   \ent{1}      E1      entregable del proyecto
%
% Al declarar cada elemento hay que poner el \label correspondiente:
%   \label{rn:02}  \label{adr:01}  \label{v:02}  \label{ca:3}
\newcommand{\rn}[1]{\hyperref[rn:#1]{\textbf{RN-#1}}}
\newcommand{\adrref}[1]{\hyperref[adr:#1]{\textsf{ADR-#1}}}
\newcommand{\vista}[1]{\hyperref[v:#1]{\textsf{V-#1}}}
\newcommand{\crit}[1]{\hyperref[ca:#1]{\textsf{CA-#1}}}
\newcommand{\ent}[1]{\textbf{E#1}}

% Identificador técnico en monoespaciado: \id{ms-reservas}, \id{usuarios_db}
\newcommand{\id}[1]{\texttt{\detokenize{#1}}}

%------------------- Bloque de decisión de arquitectura ----------------
% Un ADR compacto: los cuatro campos que hacen falta para que una decisión
% se pueda discutir. Sin «alternativa descartada» un ADR no dice nada.
%
% Uso:
%   \decision{01}{Título}{Contexto}{Decisión}{Alternativa descartada}{Consecuencia}
\newcommand{\decision}[6]{%
  \begin{tcolorbox}[enhanced,breakable,
      colback=white,colframe=acento!45,boxrule=0.6pt,arc=2pt,
      left=8pt,right=8pt,top=6pt,bottom=6pt,
      title={\textsf{ADR-#1}\quad #2},
      coltitle=white,colbacktitle=acento!85,fonttitle=\bfseries\small]
    \label{adr:#1}%
    \begin{description}[leftmargin=0pt,labelsep=6pt,itemsep=3pt]
      \item[Contexto.] #3
      \item[Decisión.] #4
      \item[Alternativa descartada.] #5
      \item[Consecuencia.] #6
    \end{description}
  \end{tcolorbox}}

%----------------- Inclusión de diagramas con aviso --------------------
% Si el archivo no existe todavía, el documento compila igual y deja un
% recuadro visible en su lugar en vez de fallar. Así nadie entrega un PDF
% al que le falta un diagrama sin enterarse.
\newcommand{\ajustar}[1]{%
  \adjustbox{max width=\linewidth,max height=0.68\textheight}{#1}}

\newcommand{\faltafig}[1]{%
  \fbox{\parbox{0.92\linewidth}{\centering\small
    \textcolor{BrickRed}{\textbf{Falta \texttt{figuras/#1.pdf}}}\\[3pt]
    Exporta el diagrama desde Structurizr y déjalo en \texttt{figuras/}
    con ese nombre.}}}

\newcommand{\diagrama}[1]{%
  \IfFileExists{figuras/#1.pdf}%
    {\ajustar{\includegraphics{#1}}}%
    {\faltafig{#1}}}

% Variante para diagramas anchos, pensada para ir dentro de un entorno
% «sidewaysfigure»: en una página girada el ancho disponible pasa a ser
% \textheight y el alto, \textwidth. Un diagrama de contenedores no cabe
% legible en vertical; girarlo casi duplica la escala.
\newcommand{\ajustarancho}[1]{%
  \adjustbox{max width=\textheight,max height=0.92\textwidth}{#1}}

\newcommand{\diagramaancho}[1]{%
  \IfFileExists{figuras/#1.pdf}%
    {\ajustarancho{\includegraphics{#1}}}%
    {\faltafig{#1}}}

% Variante «a página completa»: además de ocupar la página entera, invade
% los márgenes. Sujeta al papel en vez de al bloque de texto, un diagrama
% gana unos 20 puntos porcentuales de escala, que es la diferencia entre
% leer los nombres de los contenedores y no leerlos.
%
% El \makebox declara el ancho que cabe pero recibe contenido mayor, así
% que desborda por igual a ambos lados. Es deliberado: es la única forma de
% salir del bloque de texto sin descolocar el centrado. OJO: ese ancho
% declarado tiene que ser el del cuerpo del float donde se use, y no es el
% mismo en vertical que en horizontal. Ponerlo mal es justo lo que hace que
% la figura se vea corrida hacia un lado.
%
%   \diagramapagina{archivo}                -> vertical, dentro de «figure»
%   \diagramapagina[horizontal]{archivo}    -> girado, dentro de «sidewaysfigure»
%
% Cuál conviene depende de la PROPORCIÓN del diagrama, no de su tamaño:
% girar solo gana escala si el diagrama es más ancho que alto. Los exports
% de Structurizr suelen salir más altos que anchos, así que casi siempre
% la respuesta correcta es «vertical».
%
% Al invadir los márgenes la figura pisaría el folio, así que esa página va
% sin él con \thisfloatpagestyle{empty} (\thispagestyle no funciona dentro
% de un float).

% Margen contra el borde del papel: total de los dos lados (1 cm cada uno).
\newcommand{\margenpapel}{2cm}

% Hueco que hay que reservar para el pie de figura.
\newcommand{\huecopie}{3\baselineskip}

% VERTICAL — dentro de un «figure» normal. El cuerpo del float mide
% \linewidth de ancho, y el espacio del papel es \paperwidth × \paperheight.
\newcommand{\ajustarpaginavertical}[1]{%
  \makebox[\linewidth]{%
    \adjustbox{max width=\dimexpr\paperwidth-\margenpapel\relax,
               max height=\dimexpr\paperheight-\margenpapel-\huecopie\relax}{#1}}}

% HORIZONTAL — dentro de un «sidewaysfigure». Al girar, el cuerpo del float
% mide \textheight de ancho y las dos dimensiones del papel se intercambian.
\newcommand{\ajustarpaginahorizontal}[1]{%
  \makebox[\textheight]{%
    \adjustbox{max width=\dimexpr\paperheight-\margenpapel\relax,
               max height=\dimexpr\paperwidth-\margenpapel-\huecopie\relax}{#1}}}

% Primer argumento opcional: orientación, «vertical» (por defecto) u
% «horizontal». Segundo: se pasa tal cual a \includegraphics, sirve sobre
% todo para recortar el blanco sobrante: [trim=izq abajo der arriba,clip].
\NewDocumentCommand{\diagramapagina}{O{vertical} O{} m}{%
  \IfFileExists{figuras/#3.pdf}%
    {\ifcsname ajustarpagina#1\endcsname
       \csname ajustarpagina#1\endcsname{\includegraphics[#2]{#3}}%
     \else
       \PackageError{plantilla}{Orientacion '#1' desconocida en
         \string\diagramapagina}{Usa 'vertical' u 'horizontal'.}%
     \fi}%
    {\faltafig{#3}}}

% Pie uniforme para las tablas que continúan en la página siguiente.
\newcommand{\tablasigue}[1]{%
  \midrule
  \multicolumn{#1}{r}{\small\itshape Continúa en la siguiente página} \\
  \endfoot
  \endlastfoot
}
